\section{Evaluation}\label{sec:evaluation}

The effectiveness of the proposed forecasting framework is evaluated
through \textbf{five experiments}, each designed to examine a different
aspect of system performance. The first experiment identifies the most
suitable machine learning model for the deterministic forecasting
branch. The second compares the prediction accuracy of all four
forecasting approaches and includes statistical significance testing
together with an upper-bound performance analysis. The third experiment
investigates the stability of the models by repeating identical
prediction tasks multiple times. The fourth evaluates the reliability of
language model outputs through an automated hallucination analysis,
while the final experiment measures the overall quality of the responses
generated by each forecasting system. Unless otherwise specified, all
experiments are conducted using the same \textbf{300 holdout
predictions} collected across \textbf{60 trading dates}.

\subsection{Experiment 1: Selection of the Machine Learning Model}\label{sec:ml-selection}

The deterministic forecasting component should only be selected after
demonstrating that it performs competitively against other candidate
models under identical experimental conditions. Therefore, five
different machine learning algorithms were evaluated using the same
dataset, feature set, walk-forward validation strategy, and five-day
embargo period. Each model independently optimized its decision
threshold using only the corresponding validation data, while the
testing dataset remained completely unseen during model selection. This
controlled comparison ensures that the chosen machine learning model is
selected based on objective experimental evidence rather than prior
assumptions or implementation preference.

\begin{table}[H]
\centering
\scriptsize
\setlength{\tabcolsep}{4pt}
\resizebox{\textwidth}{!}{%
\begin{tabular}{@{}llcccccc@{}}
\toprule
\textbf{Model} & \textbf{Family} & \textbf{Overall} & \textbf{Calm} & \textbf{Shock} & \textbf{F1} & \textbf{AUC} & \textbf{Brier} \\
\midrule
Logistic Regression & Classical & 53.13\% & 53.42\% & 51.67\% & 66.74 & 51.07 & 0.2536 \\
Random Forest & Classical & 52.17\% & 52.76\% & 49.17\% & 60.21 & 50.79 & 0.2543 \\
XGBoost & Classical & 52.90\% & 53.77\% & 48.47\% & 64.56 & 50.91 & 0.2730 \\
GRU & Deep learning & 54.75\% & 54.64\% & 55.28\% & 66.62 & 50.20 & 0.2561 \\
\textbf{LSTM (selected)} & \textbf{Deep learning} & \textbf{53.84\%} & \textbf{53.96\%} & \textbf{53.19\%} & \textbf{66.15} & \textbf{50.04} & \textbf{0.2561} \\
Baseline: always UP & Baseline & 54.98\% & 55.38\% & 52.92\% & 70.95 & 50.00 & 0.2475 \\
\bottomrule
\end{tabular}}
\caption{Five candidate models on identical walk-forward folds. The selected model is highlighted.}
\label{tab:candidate-models}
\end{table}

The primary objective of this experiment is to determine whether
sequence-based deep learning models provide a meaningful advantage over
traditional machine learning algorithms, particularly during periods of
high market volatility. As shown in \textbf{Table 3}, the two sequence
models (LSTM and GRU) achieved an average shock-market accuracy of
\textbf{54.23\%}, whereas the three classical models obtained an average
accuracy of \textbf{49.77\%}. This \textbf{4.5 percentage-point
improvement} represents the largest performance difference observed in
the comparison. Furthermore, all classical models performed at or below
the \textbf{always-UP} baseline during volatile market conditions, while
both sequence models consistently exceeded this baseline. These findings
support the selection of a recurrent neural network as the deterministic
forecasting model.

Although the \textbf{GRU} achieved the highest accuracy during
shock-market periods, the performance difference between the GRU and the
LSTM was not statistically significant. The paired \textbf{McNemar test}
produced a \textbf{p-value of 0.393}, indicating that the observed
difference could have occurred by chance for the available dataset.
Consequently, the LSTM was selected as the final deterministic model
because it had already been integrated into the forecasting pipeline and
extensively validated throughout the development process. Replacing the
architecture based on such a small and statistically insignificant
improvement would increase the risk of overfitting rather than providing
a genuine performance advantage.

Another important observation from \textbf{Table 3} is that none of the
evaluated machine learning models outperformed the \textbf{always-UP
baseline} when overall accuracy was considered. In addition, the
\textbf{AUC values}, which range from \textbf{50.0} to \textbf{51.1},
indicate that all five models perform only marginally better than random
ranking. Although this result may appear disappointing, it provides an
important motivation for the remainder of this research. The limited
predictive capability of conventional price-based models suggests that
incorporating additional sources of information, such as language
reasoning and retrieval-based knowledge, may offer greater potential for
improving financial forecasting performance.

\subsection{Experiment 2: Direction Accuracy}\label{sec:direction-accuracy}

\begin{table}[H]
\centering
\scriptsize
\setlength{\tabcolsep}{4pt}
\resizebox{\textwidth}{!}{%
\begin{tabular}{@{}lcccccc@{}}
\toprule
\textbf{System} & \textbf{Overall} & \textbf{Calm} & \textbf{Shock} & \textbf{95\% CI} & \textbf{Brier (raw)} & \textbf{Brier (cal.)} \\
\midrule
\textbf{Hybrid (ML + RAG-LLM)} & \textbf{58.33\%} & \textbf{57.89\%} & \textbf{59.09\%} & \textbf{[52.33, 64.67]} & \textbf{0.2447} & \textbf{0.2447} \\
LLM + RAG (grounded) & 57.67\% & 56.32\% & 60.00\% & [51.00, 64.33] & 0.2446 & 0.2505 \\
LLM only (no evidence) & 56.00\% & 55.79\% & 56.36\% & [49.66, 62.67] & 0.2517 & 0.2509 \\
ML only (LSTM) & 52.67\% & 55.79\% & 47.27\% & [46.33, 59.33] & 0.2542 & 0.2507 \\
Baseline: always UP & 54.33\% & 54.21\% & 54.55\% & -- & -- & -- \\
\bottomrule
\end{tabular}}
\caption{Direction accuracy on 300 holdout predictions across 60 dates. The calibrated Brier column re-centres every arm on its own tuned threshold, which is what makes the calibration comparison like-for-like.}
\label{tab:direction-accuracy}
\end{table}

\begin{figure}[H]
\centering
\includegraphics[width=0.76\textwidth]{fig-accuracy-regime.png}
\caption{Accuracy by market regime against the always-UP baseline. The ML arm's collapse in shock conditions, and the language arms' resilience there, is the clearest pattern in the study.}
\label{fig:accuracy-regime}
\end{figure}

The direction prediction results for all four forecasting models are
presented in \textbf{Table 4}. Among the evaluated approaches, the
\textbf{hybrid model} achieved the highest overall accuracy of
\textbf{58.33\%}, exceeding the \textbf{always-UP baseline} by
approximately \textbf{four percentage points}. It also demonstrated the
best probability calibration, with a \textbf{Brier score of 0.2447},
compared with \textbf{0.2505} for the retrieval-enhanced language model.

A closer examination of the results across different market regimes
provides greater insight than the overall accuracy alone. During
\textbf{CALM} market conditions, the performance of all four forecasting
models is very similar, with an accuracy difference of only \textbf{2.1
percentage points}. Their results remain close to the baseline strategy,
indicating that no individual approach offers a substantial advantage in
stable market environments.

A different pattern emerges during \textbf{SHOCK} market conditions. The
LSTM model achieves an accuracy of \textbf{47.27\%}, which is
approximately \textbf{seven percentage points below} the always-UP
baseline and represents the weakest result among all evaluated models.
In contrast, the retrieval-enhanced language model reaches
\textbf{60.0\%} accuracy under the same conditions. These findings
suggest that the deterministic sequence model struggles during periods
of high market volatility, whereas the grounded language model remains
comparatively robust. This complementary behaviour provides the primary
motivation for combining both models within the proposed hybrid
framework.

Although the hybrid model achieves the highest overall accuracy, the
statistical analysis indicates that the observed improvements should be
interpreted with caution. As shown in \textbf{Table 4}, the confidence
intervals of all forecasting models overlap considerably. Furthermore,
the pairwise comparisons reported in \textbf{Table 5} show that
\textbf{none of the McNemar tests achieve statistical significance (p
\textless{} 0.05)}. In particular, the comparison between the hybrid
model and the grounded language model produces a \textbf{p-value of
0.885}, indicating that no statistically significant difference exists
between the two best-performing approaches.

Therefore, the results should not be interpreted as evidence that one
forecasting model consistently outperforms the others. Instead, the main
contribution of this experiment is the fair and systematic comparison of
four forecasting architectures under identical experimental conditions.
The analysis demonstrates the strengths and limitations of each
engineering component, providing insight into the contribution of
language reasoning, retrieval grounding, and model fusion beyond the
overall prediction accuracy.

\begin{table}[H]
\centering
\scriptsize
\setlength{\tabcolsep}{4pt}
\resizebox{\textwidth}{!}{%
\begin{tabular}{@{}lccccc@{}}
\toprule
\textbf{Comparison (A vs B)} & \textbf{Rows} & \textbf{A right, B wrong} & \textbf{B right, A wrong} & \textbf{p (exact)} & \textbf{p $<$ 0.05?} \\
\midrule
ML vs LLM & overall & 53 & 63 & 0.4035 & no \\
ML vs LLM & shock & 17 & 27 & 0.1742 & no \\
ML vs LLM+RAG & overall & 60 & 75 & 0.2281 & no \\
ML vs LLM+RAG & shock & 21 & 35 & 0.0814 & no \\
ML vs Hybrid & overall & 35 & 52 & 0.0857 & no \\
ML vs Hybrid & shock & 17 & 30 & 0.0789 & no \\
LLM vs Hybrid & overall & 52 & 59 & 0.5692 & no \\
LLM+RAG vs Hybrid & overall & 23 & 25 & 0.8854 & no \\
\bottomrule
\end{tabular}}
\caption{Exact paired McNemar tests. Only rows on which two systems disagree carry information, which is why this is the right test for paired predictions.}
\label{tab:mcnemar-tests}
\end{table}

The strongest signal in the study is the grounded language arm against
the ML arm in shock markets, at p = 0.081, closely followed by fusion
against ML in shock at p = 0.079. Neither clears 5\%, and both are
reported as directional rather than proven. With 110 shock rows this is
the resolution the design can deliver; Section 7 returns to what would
be needed to settle it.

Because the accuracies sit close to the baselines, the natural question
is whether more engineering would lift them. Table 6 answers it. The
last three rows are deliberately unfair - they are computed by fitting
on the test set - and they bound what any amount of tuning on these four
arms could achieve.

\begin{table}[H]
\centering
\footnotesize
\begin{tabularx}{\textwidth}{@{}Ycp{4.6cm}@{}}
\toprule
\textbf{Bound on the same 300 holdout rows} & \textbf{Accuracy} & \textbf{Legitimate?} \\
\midrule
Always predict DOWN & 45.67\% & baseline \\
Always predict UP & 54.33\% & baseline \\
\textbf{Reported hybrid (weights and thresholds from the tuning split)} & \textbf{58.33\%} & \textbf{yes -- this is what is reported} \\
Each arm's threshold refitted on the holdout itself & 58.67\% & no -- fitted on the test set \\
Best three-arm blend, weights and cut-off fitted on the holdout & 60.33\% & no -- fitted on the test set \\
Oracle: per row, believe whichever arm is right & 82.67\% & impossible -- needs the answer first \\
\bottomrule
\end{tabularx}
\caption{What is attainable on this task. Refitting every threshold on the test set buys 0.34 points over the honestly reported figure.}
\label{tab:attainable-bounds}
\end{table}

Two things follow. First, refitting every threshold on the test set - a
serious form of cheating - buys 0.34 percentage points, which says the
reported result is not sensitive to threshold choice and that the honest
and dishonest numbers are nearly the same here. Second, on 17.33\% of
holdout rows all four arms are wrong simultaneously, so even the oracle
tops out at 82.67\%. Any reported accuracy above that on this task would
indicate a look-ahead leak or test-set tuning rather than a better
model.

\subsection{Experiment 3: Stability Under Identical Input}\label{sec:stability}

A deterministic model and a non-deterministic one cannot be compared on
a single run. The identical prompt was therefore sent five times on each
of 30 cases, at temperature 0.

\begin{table}[H]
\centering
\scriptsize
\setlength{\tabcolsep}{4pt}
\resizebox{\textwidth}{!}{%
\begin{tabular}{@{}lcccc@{}}
\toprule
\textbf{System} & \textbf{Answer changed on identical input} & \textbf{Mean sd of p(up)} & \textbf{Mean agreement} & \textbf{Notes} \\
\midrule
ML only (LSTM) & 0.00\% & 0.0000 & 1.0000 & deterministic by construction \\
LLM + RAG & 3.33\% & 0.0079 & 0.9867 & shock 6.67\%, calm 0.00\% \\
Hybrid & 3.33\% & 0.0040 & 0.9867 & moves 0.505$\times$ as much as the LLM \\
\bottomrule
\end{tabular}}
\caption{Repeated identical prompts, 5 runs on each of 30 cases.}
\label{tab:stability}
\end{table}

\begin{figure}[H]
\centering
\includegraphics[width=0.70\textwidth]{fig-stability.png}
\caption{Spread of each arm's probability across five identical prompts.}
\label{fig:stability}
\end{figure}

\textbf{Fusion halves the movement but does not confer determinism.} The
hybrid\textquotesingle s probability moves 0.505 times as much as the
language model\textquotesingle s across identical prompts, which is
close to the weight sitting on the deterministic half and is what the
arithmetic predicts. Its direction flip rate, however, is identical at
3.33\%, and it flipped on exactly the cases where the language model
flipped and on no others. Blending damps the size of a disagreement, not
the fact of one: halving the movement does not help when a case already
sits on the decision boundary, which is precisely where a flip happens.

The hypothesis that the hybrid would inherit the ML
arm\textquotesingle s stability is therefore supported on variance and
not supported on determinism. That distinction was only visible because
it was measured rather than inferred from the fact that half the hybrid
is a fixed function, and it is a caution for any system that reports a
fused prediction as though it were reproducible. All the instability
sits in shock markets, at 6.67\% against 0.00\% in calm conditions - the
same regime where the language arms are most useful.

\subsection{Experiment 4: Hallucination Audit}\label{sec:hallucination-audit}

\begin{table}[H]
\centering
\footnotesize
\begin{tabularx}{\textwidth}{@{}YCC@{}}
\toprule
\textbf{Check} & \textbf{LLM only} & \textbf{LLM + RAG} \\
\midrule
Fabricated a figure it was not shown & 0.0\% & 0.0\% \\
\rowcolor{yellow!18}\textbf{Unit error (right number, wrong scale)} & \textbf{2.6\%} & \textbf{14.2\%} \\
Invented a citation & 0.0\% & 0.0\% \\
Named the real company despite anonymisation & 0.0\% & 0.0\% \\
Cited the evidence it was given & n/a -- shown none & 38.4\% \\
Fully clean on all checks & 97.4\% & 85.8\% \\
Responses audited & 500 & 500 \\
\bottomrule
\end{tabularx}
\caption{Three mechanical checks applied to every response. No human and no language model rated anything.}
\label{tab:hallucination-audit}
\end{table}

\begin{figure}[H]
\centering
\includegraphics[width=0.67\textwidth]{fig-hallucination.png}
\caption{The mechanical hallucination checks by arm.}
\label{fig:hallucination}
\end{figure}

This experiment did not produce the result it was designed to find. The
premise - taken from the retrieval literature - was that an ungrounded
model would fabricate figures and that retrieval would reduce the rate.
Neither arm fabricated a single figure across 1,000 responses, and
neither invented a single citation, so there was essentially no
hallucination for retrieval to remove. The anonymisation control also
held completely: the model never named the real company, which means the
accuracy figures in Section 6.2 cannot be explained by the model
recalling what these five stocks actually did.

\textbf{The one failure that did occur got worse with retrieval.} Unit
errors - restating a return of -14.57\% as -0.1457, keeping the
magnitude and losing the scale - rose from 2.6\% to 14.2\%, a factor of
five. More context meant more numbers to keep straight. This is a real
cost of grounding and it is not one the retrieval literature prepares
you for, because the standard evaluation looks for invented propositions
rather than mis-scaled figures. For a trading assistant the distinction
matters: a user who acts on -0.1457 believing it to be a percentage has
been misled just as effectively as by a fabrication.

The grounded arm cited at least one retrieved item on 38.4\% of
responses. The remaining responses judged on the price data alone, which
the prompt permits when the evidence does not bear on the question.
Retrieval, then, changed the answer on roughly two in five cases and was
correctly ignored on the rest.

\subsection{Experiment 5: Response Quality as a Trading Assistant}\label{sec:response-quality}

Direction accuracy asks whether the call was right, and Section 6.2
showed that on this task that question is barely answerable: six points
separate the four arms and the confidence intervals all overlap. But a
user does not consume a call. They consume an answer - a direction, a
reason, evidence they can check, and advice shaped to their risk
appetite. Two systems can be equally accurate and differ enormously in
whether that answer can be trusted or acted on, and for a trading
assistant a confident, unverifiable, unhedged answer is the expensive
failure rather than a coin flip that landed badly.

\begin{figure}[H]
\centering
\includegraphics[width=0.94\textwidth]{fig-response-quality.png}
\caption{Response quality across the four rubric dimensions, scored mechanically over the same 300 holdout cases. This is the axis on which the four systems actually separate: 26.6 points, against 5.66 points on accuracy.}
\label{fig:response-quality}
\end{figure}

\begin{table}[H]
\centering
\scriptsize
\setlength{\tabcolsep}{4pt}
\resizebox{\textwidth}{!}{%
\begin{tabular}{@{}lccccc@{}}
\toprule
\textbf{System} & \textbf{Grounding /30} & \textbf{Personalisation /25} & \textbf{Helpfulness /25} & \textbf{Safety /20} & \textbf{Total /100} \\
\midrule
\rowcolor{yellow!18}\textbf{Hybrid (ML + RAG-LLM)} & \textbf{24.80} & \textbf{16.95} & \textbf{22.94} & \textbf{13.33} & \textbf{78.0} \\
LLM + RAG (grounded) & 24.80 & 18.47 & 22.94 & 11.00 & 77.2 \\
LLM only (no evidence) & 22.90 & 17.11 & 18.88 & 11.00 & 69.9 \\
ML only (LSTM) & 23.00 & 10.39 & 7.00 & 11.00 & 51.4 \\
\bottomrule
\end{tabular}}
\caption{The 17-rule rubric across four dimensions, applied to 300 responses per arm.}
\label{tab:response-quality}
\end{table}

This is the experiment that separates the four systems. On accuracy they
span 5.66 points and the difference is not significant; on response
quality they span 26.6 points, from 51.4 for the LSTM to 78.0 for the
fusion arm. The ordering is also the ordering of engineering effort,
which is not true of the accuracy table.

\textbf{How to read the near-tie at the top.} The fusion arm scores
highest at 78.0, but only 0.8 points clear of the grounded arm. That
near-tie is not noise - it is structural, and it is the finding. On
grounding and helpfulness the two are identical to two decimal places,
because the fusion arm\textquotesingle s answer text is the grounded
arm\textquotesingle s answer text: it inherits the reasoning and the
evidence wholesale and cannot differ. The fusion gains on safety, where
it is better calibrated and is the only arm able to tell a user that its
two components disagreed. It loses on personalisation, because pooling
concentrates its confidence into a narrower band, so the risk profiles
separate its advice less often.

The conclusion that follows is specific and, I think, the most useful
thing in this report: retrieval is where response quality comes from,
and fusion adds calibration and disagreement-reporting worth about a
point. Grounding lifts the language model from 69.9 to 77.2 - a
7.3-point gain, nine times the gain from fusion. The LSTM scores 51.4
not because it is inaccurate, but because a bare probability cannot
explain itself, attribute a claim, or be checked; it takes 7.0 of 25 on
helpfulness and 10.39 of 25 on personalisation.

\textbf{Two scoring decisions that changed the answer.} First,
text-quality rules are applied only to model-generated text and never to
the deterministic fusion sentence. An earlier version scored the whole
response, and the fusion arm took full marks for situational awareness
purely because that template contains the word ``market'' - which graded
this repository\textquotesingle s prose rather than the system. It
scored 86.2 that way against 78.0 here, and the earlier figure was
wrong. Second, a silent system is not penalised for what it cannot say:
the LSTM takes full marks for never fabricating a figure, having
produced no text, and loses points instead on attribution, explanation
and evidence, which is where the absence actually costs a user.

One limitation of the rubric should be stated. Of the 100 points, 42.0
are earned identically by all four arms - rules such as never
fabricating a figure and never using absolutist language. Those rules
are worth keeping, and reporting that no arm ever violated them is
itself a result, but they inflate all four totals equally. The
differences between systems, not the totals, are what carry information.

\subsection{Discussion}\label{sec:discussion}

An upper-bound analysis was performed to estimate the maximum achievable
accuracy for the task.

The oracle model achieved an accuracy of approximately 82.67\%.

This result demonstrates that further improvements are theoretically
possible while also illustrating the limitations imposed by the
available data.
